\section{State}
\label{sec:pat-state}

\problemstatement{Allow an object to alter its behaviour when its internal
state changes, so it appears to change its class -- replacing sprawling
conditionals on a state variable with a set of state objects.}

\begin{patternbox}[When You Reach for It]
The trigger is \emph{``behaviour depends on a mode/state, and the same
method does different things in each''} -- a vending machine, a media
player (playing/paused/stopped), a TCP connection, an order lifecycle.
When methods are riddled with \texttt{if (state == ...)}, extract states
into classes.
\end{patternbox}

\subsection*{Understanding the pattern}

Each state becomes a class implementing a common \texttt{State} interface;
the context holds a pointer to the current state and delegates behaviour to
it. A state's methods perform the state-specific action and can transition
the context to another state. Adding a new state is a new class, not
another branch in every method -- Open/Closed for state machines.

\begin{ideabox}[One Class per State, Delegate and Transition]
Turn each mode into an object that knows how to behave \emph{and} which
state comes next. The context just forwards calls to its current state
object. The scattered \texttt{switch(state)} logic collapses into
localised, self-contained state classes.
\end{ideabox}

\subsection*{C++ implementation}

\begin{lstlisting}
#include <bits/stdc++.h>
using namespace std;

class Fan;
struct State {
    virtual void pull(Fan& fan) = 0;            // the chain-pull action
    virtual string name() const = 0;
    virtual ~State() = default;
};

class Fan {
    unique_ptr<State> state;
public:
    Fan(unique_ptr<State> s) : state(move(s)) {}
    void setState(unique_ptr<State> s) { state = move(s); }
    void pull() { state->pull(*this); }
    string status() const { return state->name(); }
};

struct Off; struct Low;
struct Low : State {
    void pull(Fan& fan) override;
    string name() const override { return "Low"; }
};
struct Off : State {
    void pull(Fan& fan) override { fan.setState(make_unique<Low>()); }
    string name() const override { return "Off"; }
};
void Low::pull(Fan& fan) { fan.setState(make_unique<Off>()); }

int main() {
    Fan fan(make_unique<Off>());
    cout << fan.status() << "\n";               // Off
    fan.pull(); cout << fan.status() << "\n";   // Low
    fan.pull(); cout << fan.status() << "\n";   // Off
    return 0;
}
\end{lstlisting}

\begin{complexitybox}[Trade-offs]
\textbf{Pros.} Eliminates large state conditionals; each state is
isolated and testable; transitions are explicit; new states are added
without editing others. \textbf{Cons.} More classes; transition logic is
spread across state classes, so the overall state machine can be harder to
see at a glance (a transition table can complement it).
\end{complexitybox}

\begin{takeawaybox}
State replaces ``behaviour depends on a mode'' conditionals with one class
per state that both acts and decides the next state. Structurally it is
Strategy, but the intent differs: State objects \emph{swap themselves} as
the machine transitions, whereas a Strategy is chosen by the client.
\end{takeawaybox}
